% Context from chapters-tex/fourier.tex; for mathematical reference, not compilation.
\section{Convolution on \texorpdfstring{$\RR$ and $\TT$}{the Line and the Circle}}

                                                
\subsection{Convolution}

\begin{figure}[tbp]
\centering
\BookDrawing{tikz/fourier-convolution}
\caption{Convolution on $\RR$. The dashed curve is $f\star g$; the marked value $(f\star g)(x)$ equals the shaded integral.}
\label{fig-review-fourier-convolution}
\end{figure}
On $\XX=\RR$ or $\TT$, one defines 
\eql{\label{eq-convol-1d}
	f \star g(x) = \int_\XX f(t) g(x-t) \d t. 
}
Here and throughout this section, integrals use Lebesgue measure on $\RR$ and normalized Haar measure $\d t/(2\pi)$ on $\TT$. With this convention, Young's inequality states that, for $1\leq p,q,r\leq\infty$ and $(f,g)\in L^p(\XX)\times L^q(\XX)$,
\eql{\label{eq-young-convol}
	\frac{1}{p}+\frac{1}{q}=1+\frac{1}{r} \qarrq
	f \star g \in L^r(\XX) \qandq 
	\norm{f\star g}_{L^r(\XX)} \leq \norm{f}_{L^p(\XX)}\norm{g}_{L^q(\XX)}.
}
In particular, convolution by $f\in L^1(\XX)$ defines a bounded linear map from $L^p(\XX)$ to itself. When $r=\infty$ and $1<p,q<\infty$, the convolution $f\star g$ has a bounded continuous representative, illustrating its regularizing effect.
 
On the circle, $p<q \qarrq L^q(\TT) \subset L^p(\TT)$, so $L^\infty(\XX)$ is contained in every space in this scale.


\myfigure{
\BookDrawing{tikz/fourier-running-average}
}{ 
	Signal filtering with a box filter (running average).   
}{fig-filter-box}

Convolution can regularize functions. For instance, if $f \in L^1(\XX)$ and $g\in C^1(\XX)$ is bounded with bounded derivative, then $f \star g$ is differentiable and $(f \star g)'=f \star g'$. To construct an approximate identity, choose a smooth compactly supported $\rho\geq0$ with $\int_\RR\rho=1$ and set $\rho_\epsilon(x)=\epsilon^{-1}\rho(x/\epsilon)$. Then $f\star\rho_\epsilon\to f$ in $L^p(\RR)$ for $1\leq p<\infty$, and uniformly for bounded uniformly continuous $f$. On the circle, use $2\pi\sum_{k\in\ZZ}\rho_\epsilon(\cdot-2\pi k)$, whose integral against normalized Haar measure is one.
 
This smoothing effect is also useful for denoising signals and images.


\myfigure{
\image{fourier}{.6}{filtering-1d-1}\\
\image{fourier}{.6}{filtering-1d-2}\\
\image{fourier}{.6}{filtering-1d-3}
}{ 
	Filtering an irregular signal with Gaussian filters of increasing width $\si$.  
}{fig-filtering-1d}

   
\begin{prop}
If $f,g\in L^1(\XX)$, then
\eql{\label{eq-convol-fourier}
	\Ff(f \star g) = \hat f \odot \hat g.
}
For $\XX=\RR$, $\odot$ denotes pointwise multiplication of the continuous Fourier transforms.
For $\XX=\TT$, $\odot$ acts on bounded Fourier-coefficient sequences (pointwise multiplication of sequences).
\end{prop}


           
\begin{figure}[tbp]
\centering
\BookDrawing{tikz/fourier-convolution-fourier}
\caption{How the Fourier transform converts convolution into multiplication.}
\label{fig-review-fourier-convolution-fourier}
\end{figure}
This means that $\Ff$ is an algebra homomorphism. For instance, if $\XX=\RR$, its range lies in the algebra of continuous functions that vanish at $\pm\infty$.

As shown in Figure~\ref{fig-splines}, successive convolutions of a box function produce cardinal splines: piecewise polynomials of increasing smoothness.

\begin{figure}[tbp]
\centering
\BookDrawing{tikz/fourier-splines}
\caption{Cardinal splines are defined by successive convolutions.}
\label{fig-splines}
\end{figure}


Convolution also describes the sum of independent random variables. Writing $f_X$ for the probability density of a random vector $X$ with respect to Lebesgue measure, independence of $X$ and $Y$ gives $f_{X+Y}=f_X \star f_Y$. This identity is fundamental in probability and statistics.


\texttt{Associated code: fourier/test\_denoising.m and fourier/test\_fft2.m}

                                                
\subsection{Translation Invariant Operators}

\begin{figure}[tbp]
\centering
\BookDrawing{tikz/fourier-translation-inv}
\caption{Commutative diagram for translation invariance.}
\label{fig-review-fourier-translation-inv}
\end{figure}
Translation-invariant operators commute with translations. In signal and image processing, they apply the same rule at every location.

The following propositions characterize translation-invariant operators as convolutions against kernels that may be distributions. The term ``translation equivariant'' is more precise: translating the input translates the output by the same amount.
 
The kernel\textquotesingle s regularity depends on the topologies of the input and output spaces, so the proof differs between the two settings below.
 
We begin with operators whose outputs are continuous functions.

\begin{prop}
	We define $T_\tau f = f(\cdot-\tau)$.
	 
	A bounded linear operator $H : (L^2(\XX),\norm{\cdot}_2) \rightarrow (C_b(\XX),\norm{\cdot}_\infty)$ commutes with every translation, meaning that for all $\tau$, $H \circ T_\tau = T_\tau \circ H$, if and 